\documentclass[12pt,a4paper]{report}

% ── Packages ──────────────────────────────────────────────────────
\usepackage[a4paper, margin=1in]{geometry}
\usepackage{times}
\usepackage{graphicx}
\usepackage{amsmath, amssymb}
\usepackage{booktabs}
\usepackage{array}
\usepackage{longtable}
\usepackage{hyperref}
\usepackage{url}
\usepackage{listings}
\usepackage{xcolor}
\usepackage{fancyhdr}
\usepackage{setspace}
\usepackage{titlesec}
\usepackage{caption}
\usepackage{subcaption}
\usepackage{float}
\usepackage{enumitem}
\usepackage{multirow}
\usepackage{algorithm}
\usepackage{algpseudocode}
\usepackage{tikz}
\usepackage{pgfplots}
\pgfplotsset{compat=1.18}
\usetikzlibrary{shapes.geometric, arrows, positioning, fit, backgrounds}

% ── Fix paragraph indentation (no indent, small skip between paras) ──
\usepackage{parskip}
\setlength{\parindent}{0pt}
\setlength{\parskip}{6pt plus 2pt minus 1pt}

% ── Colors ────────────────────────────────────────────────────────
\definecolor{codegray}{RGB}{245, 245, 245}
\definecolor{codegreen}{RGB}{0, 128, 0}
\definecolor{codepurple}{RGB}{128, 0, 128}
\definecolor{codeblue}{RGB}{0, 82, 147}
\definecolor{winnergreen}{RGB}{220, 255, 220}
\definecolor{lightblue}{RGB}{235, 245, 255}

% ── Hyperref setup ────────────────────────────────────────────────
\hypersetup{
    colorlinks=true,
    linkcolor=black,
    citecolor=black,
    urlcolor=black,
}

% ── Code listing style ────────────────────────────────────────────
\lstdefinestyle{pythonstyle}{
    backgroundcolor=\color{codegray},
    commentstyle=\color{codegreen},
    keywordstyle=\color{codeblue}\bfseries,
    stringstyle=\color{codepurple},
    basicstyle=\ttfamily\footnotesize,
    breaklines=true,
    frame=single,
    language=Python,
    showstringspaces=false,
    tabsize=4,
    captionpos=b
}

% ── Page style ────────────────────────────────────────────────────
\pagestyle{fancy}
\fancyhf{}
\fancyhead[L]{\small Lightweight Session-Based Recommendation System}
\fancyhead[R]{\small UPES, Dehradun}
\fancyfoot[C]{\thepage}
\renewcommand{\headrulewidth}{0.4pt}

% ── Chapter/Section title format ──────────────────────────────────
\titleformat{\chapter}[block]
  {\normalfont\large\bfseries\centering}
  {\thechapter.}{0.5em}{\large\bfseries\MakeUppercase}
\titlespacing*{\chapter}{0pt}{12pt}{12pt}

\titleformat{\section}[block]
  {\normalfont\normalsize\bfseries}
  {\thesection.}{0.5em}{}
\titlespacing*{\section}{0pt}{10pt}{6pt}

\titleformat{\subsection}[block]
  {\normalfont\normalsize\bfseries}
  {\thesubsection.}{0.5em}{}
\titlespacing*{\subsection}{0pt}{8pt}{4pt}

\titleformat{\subsubsection}[block]
  {\normalfont\normalsize\bfseries}
  {\thesubsubsection.}{0.5em}{}
\titlespacing*{\subsubsection}{0pt}{6pt}{4pt}

% ── Line spacing ──────────────────────────────────────────────────
\onehalfspacing

\setcounter{secnumdepth}{4}
\setcounter{tocdepth}{4}

% ─────────────────────────────────────────────────────────────────
\begin{document}

% ═══════════════════════════════════════════════════════════════
% TITLE PAGE
% ═══════════════════════════════════════════════════════════════
\begin{titlepage}
\centering
\vspace*{1cm}

{\Large \textbf{Lightweight Session-Based Recommendation System}}\\[0.5cm]
{\large \textit{A}}\\[0.3cm]
{\large \textit{\textbf{Project Report}}}\\[0.2cm]
{\normalsize \textit{submitted in partial fulfillment of the}}\\
{\normalsize \textit{requirements for the award of the degree of}}\\[0.5cm]

{\Large \textbf{BACHELOR OF TECHNOLOGY}}\\[0.2cm]
{\large \textbf{in}}\\[0.2cm]
{\Large \textbf{COMPUTER SCIENCE \& ENGINEERING}}\\[0.8cm]

{\normalsize \textbf{by}}\\[0.5cm]

\begin{tabular}{ll}
\textbf{Name} & \textbf{Roll No.} \\
Saksham Agrawal & R2142231383 \\
Abhishek Yadav  & R2142231410 \\
Yash Dagar      & R2142231424 \\
\end{tabular}\\[0.8cm]

{\normalsize \textit{under the guidance of}}\\[0.3cm]
{\large \textbf{Dr. Amar Jindal}}\\[1.5cm]

\includegraphics[width=4cm]{upes_logo.png}\\[0.5cm]

{\large \textbf{School of Computer Science}}\\
{\large \textbf{University of Petroleum \& Energy Studies}}\\
{\normalsize Bidholi, Via Prem Nagar, Dehradun, Uttarakhand}\\[0.4cm]
{\normalsize \textbf{January -- May, 2026}}

\end{titlepage}

% ── Roman numeral page numbering for front matter ─────────────────
\pagenumbering{roman}

% ═══════════════════════════════════════════════════════════════
% CANDIDATE'S DECLARATION
% ═══════════════════════════════════════════════════════════════
\chapter*{Candidate's Declaration}
\addcontentsline{toc}{chapter}{Candidate's Declaration}

I/We hereby certify that the project work entitled
\textbf{``Lightweight Session-Based Recommendation System''}
in partial fulfillment of the requirements for the award of the Degree of
\textbf{BACHELOR OF TECHNOLOGY} in
\textbf{COMPUTER SCIENCE AND ENGINEERING}
with specialization in \textbf{Artificial Intelligence and Machine Learning}
and submitted to the Department of Artificial Intelligence, School of Computer Science,
University of Petroleum \& Energy Studies, Dehradun, is an authentic record of
my/our work carried out during a period from \textbf{January, 2026} to
\textbf{May, 2026} under the supervision of
\textbf{Dr.\ Amar Jindal, Assistant Professor, Data Science Cluster, UPES}.

The matter presented in this project has not been submitted by me/us for the award
of any other degree of this or any other University.

\vspace{2cm}

\noindent
\begin{tabular}{p{4cm} p{4cm} p{4cm}}
\textbf{(Saksham Agrawal)} & \textbf{(Abhishek Yadav)} & \textbf{(Yash Dagar)} \\
\textbf{Roll No.\ R2142231383} & \textbf{Roll No.\ R2142231410} & \textbf{Roll No.\ R2142231424} \\
\end{tabular}

\vspace{1.5cm}
\noindent
This is to certify that the above statement made by the candidate is correct to the
best of my knowledge.

\vspace{1.5cm}

\noindent
Date: \underline{\hspace{3cm}} \hfill \textbf{Dr.\ Amar Jindal}\\
\hfill Project Guide

% ═══════════════════════════════════════════════════════════════
% ACKNOWLEDGEMENT
% ═══════════════════════════════════════════════════════════════
\chapter*{Acknowledgement}
\addcontentsline{toc}{chapter}{Acknowledgement}

We wish to express our deep gratitude to our guide \textbf{Dr.\ Amar Jindal},
for all advice, encouragement and constant support he has given us throughout our
project work. This work would not have been possible without his support and
valuable suggestions on integrating Kolmogorov-Arnold Networks into sequential
recommendation systems.

We sincerely thanks to our respected Head of Department, School of Computer
Science, UPES, for great support in doing our project, and for providing access
to the High-Performance Computing (HPC) cluster which was essential for training
our deep learning models.

We are also grateful to Dean SoCS, UPES for giving us the necessary facilities
to carry out our project work successfully. We also thanks to our Course
Coordinator and Activity Coordinator for providing timely support and information
during the completion of this project.

We acknowledge the use of the \textbf{YooChoose} dataset provided by the RecSys 2015
challenge organizers, which served as the benchmark for evaluating all our proposed
architectures.

We would like to thank all our friends for their help and constructive criticism
during our project work. Finally, we have no words to express our sincere gratitude
to our parents who have shown us this world and for every support they have given us.

\vspace{1.5cm}

\noindent
\begin{tabular}{p{4cm} p{4cm} p{4cm}}
\textbf{Saksham Agrawal} & \textbf{Abhishek Yadav} & \textbf{Yash Dagar} \\
R2142231383 & R2142231410 & R2142231424 \\
\end{tabular}

% ═══════════════════════════════════════════════════════════════
% ABSTRACT
% ═══════════════════════════════════════════════════════════════
\chapter*{Abstract}
\addcontentsline{toc}{chapter}{Abstract}

Session-based recommendation systems have emerged as a critical solution for
modern digital platforms where user identity is often anonymous, transient, or
privacy-restricted. Unlike traditional recommendation systems that rely on
long-term user profiles, session-based approaches predict the next item a user
will interact with by analyzing only the current click sequence within a single
session.

This project proposes and evaluates a family of architectures that integrate
\textbf{Kolmogorov-Arnold Networks (KANs)} with recurrent neural network
frameworks to address the limitations of traditional RNN-MLP-based recommenders.
Three architectures are designed, implemented, and systematically compared on the
\textbf{YooChoose} e-commerce dataset (50,000 sessions, 11,910 unique items):

\begin{enumerate}
    \item \textbf{RNN Baseline}: Standard 2-layer GRU encoder with embedding and
          fully connected output layer.
    \item \textbf{Pure KAN}: GRU entirely replaced by a recurrent KAN layer ---
          each time step computes $h_t = \text{KAN}(\text{concat}(x_t, h_{t-1}))$.
    \item \textbf{Hybrid (GRU + KAN)}: Two-level recurrence --- GRU processes raw
          embeddings first, then a recurrent KAN layer further refines the GRU outputs.
\end{enumerate}

All models use B-spline parameterized learnable activation functions with grid
size 8 and spline order 3. Evaluation metrics include Top-1 Accuracy, Hit Ratio
at 10 (HR@10), Hit Ratio at 20 (HR@20), Normalized Discounted Cumulative Gain
at 10 (NDCG@10), Mean Reciprocal Rank (MRR@20), training time, and parameter
count. Experimental results on 128,029 total samples show the RNN baseline
achieves the best predictive performance (HR@10 = 0.5808, NDCG@10 = 0.3942) in
substantially less training time (145.8 s), while KAN-based architectures
demonstrate enhanced interpretability through learnable B-spline transformations
at the cost of additional parameters and compute.

\vspace{0.3cm}
\noindent\textbf{Keywords:} Session-based Recommendation, Kolmogorov-Arnold Networks,
B-Splines, Recurrent Neural Networks, GRU, YooChoose, Deep Learning

% ═══════════════════════════════════════════════════════════════
% TABLE OF CONTENTS / FIGURES / TABLES
% ═══════════════════════════════════════════════════════════════
\tableofcontents
\listoffigures
\listoftables

% ── Switch to Arabic page numbering for main content ─────────────
\pagenumbering{arabic}

% ═══════════════════════════════════════════════════════════════
% CHAPTER 1: INTRODUCTION
% ═══════════════════════════════════════════════════════════════
\chapter{Introduction}

\section{Background and Context}

The exponential growth of digital platforms has transformed how users discover
and consume content. Recommendation systems have become indispensable tools for
enhancing user experience and driving engagement on e-commerce sites, streaming
services, and news portals. These systems analyze user behavior to predict
preferences and suggest relevant items, addressing the information overload
problem in modern digital ecosystems.

Traditional approaches such as collaborative filtering rely heavily on persistent
user profiles built from long-term historical data. These methods face significant
challenges in contexts where user identity is unknown, sessions are short-lived,
or privacy regulations restrict tracking. The rise of privacy-conscious users,
anonymous browsing, and regulations like GDPR has amplified the need for
recommendation systems that operate effectively without long-term user profiles.

\section{Session-Based Recommendation}

\textbf{Session-Based Recommendation Systems} represent a paradigm shift in
recommendation technology. These systems predict the next item a user will
interact with by analyzing only the current interaction sequence within a single
session, without relying on historical user profiles. A session typically consists
of a sequence of user actions --- clicks, views, or searches --- occurring within
a limited timeframe.

Formally, given a session $S = \{i_1, i_2, \ldots, i_{n-1}\}$, the task is to
predict item $i_n$. The challenge lies in extracting meaningful patterns from
short, noisy, and highly variable sequences to make accurate predictions.

\section{Kolmogorov-Arnold Networks}

\textbf{Kolmogorov-Arnold Networks (KANs)}, introduced by Liu et al.\ (2024),
are a novel neural architecture based on the Kolmogorov-Arnold representation
theorem. The theorem states that any multivariate continuous function
$f(x_1, \ldots, x_n)$ can be written as:

\begin{equation}
f(x_1, \ldots, x_n) = \sum_{q=0}^{2n} \Phi_q \left( \sum_{p=1}^{n} \phi_{q,p}(x_p) \right)
\end{equation}

where $\phi_{q,p}: [0,1] \to \mathbb{R}$ and $\Phi_q: \mathbb{R} \to \mathbb{R}$
are learnable univariate functions.

Unlike traditional MLPs that use fixed activation functions (ReLU, sigmoid, tanh)
applied uniformly across neurons, KANs employ \textbf{learnable univariate
B-spline functions} on every network edge. This enables:

\begin{itemize}
    \item \textbf{Superior function approximation}: KANs represent complex
          functions more efficiently than MLPs.
    \item \textbf{Enhanced interpretability}: Learned B-spline functions provide
          direct insights into feature transformations.
    \item \textbf{Parameter efficiency}: Reduced redundancy compared to MLP
          architectures.
    \item \textbf{Mathematical grounding}: Based on rigorous theoretical
          foundations.
\end{itemize}

\section{Problem Statement}

Design a session-based recommendation system that predicts the next most probable
item from short, noisy interaction sequences without relying on long-term user
profiles. The solution must achieve:

\begin{itemize}
    \item \textbf{High prediction accuracy}: Competitive performance on standard
          benchmarks.
    \item \textbf{Computational efficiency}: Practical parameter usage.
    \item \textbf{Interpretability}: Transparent feature transformations through
          learnable B-splines.
    \item \textbf{Robustness}: Effective handling of sparse and short sequences.
\end{itemize}

\section{Project Objectives}

\begin{enumerate}
    \item Implement three architectures: RNN baseline, Pure Recurrent KAN, and
          Hybrid GRU+KAN.
    \item Design a custom recurrent KAN cell where
          $h_t = \text{KAN}(\text{concat}(x_t, h_{t-1}))$.
    \item Evaluate all models on the YooChoose dataset using HR@10, HR@20,
          NDCG@10, NDCG@20, MRR@20, and Top-1 Accuracy.
    \item Compare parameter efficiency, training time, and predictive performance
          across all three architectures.
\end{enumerate}

\section{Report Organization}

The remainder of this report is organized as follows: Chapter 2 reviews related
literature. Chapter 3 describes the dataset and preprocessing. Chapter 4 details
all three model architectures. Chapter 5 covers the experimental setup and
training procedure. Chapter 6 presents results and discussion. Chapter 7 concludes
the work.

% ═══════════════════════════════════════════════════════════════
% CHAPTER 2: LITERATURE REVIEW
% ═══════════════════════════════════════════════════════════════
\chapter{Literature Review}

\section{Session-Based Recommendation Systems}

The evolution of session-based recommendation systems has been driven by the need
to handle scenarios where traditional collaborative filtering fails. Early methods
relied on item-to-item similarity and sequential pattern mining, which struggled
to capture complex temporal dynamics.

\textbf{Hidasi et al.\ (2015)} \cite{hidasi2015} introduced a groundbreaking
approach using RNNs with GRU and LSTM architectures for session-based
recommendations. Their work on the YooChoose dataset demonstrated that RNN-based
methods significantly outperform standard non-sequential and neighborhood methods.
However, their approach revealed limitations including dependency on long sequences
and challenges with data sparsity.

\textbf{Li et al.\ (2017)} \cite{li2017} advanced the field with NARM (Neural
Attentive Recommendation Machine), which combines GRU encoders with attention
mechanisms. By modeling both sequential behavior and the user's main purpose
within a session, NARM achieved state-of-the-art performance. Despite these
improvements, the model remained a black-box without inherent visualization of
learned features.

\textbf{Sun et al.\ (2019)} \cite{sun2019} introduced BERT4Rec, applying
bidirectional Transformer encoder representations to sequential recommendation.
While achieving impressive performance, Transformer-based models typically require
significantly more parameters --- a consideration that motivates the development
of lightweight KAN-RNN alternatives.

\textbf{Kang and McAuley (2018)} \cite{kang2018} proposed SASRec, a
self-attentive sequential recommendation model based on unidirectional
Transformers. Their work demonstrated that attention mechanisms can capture
long-range item dependencies more effectively than RNNs alone.

\section{Kolmogorov-Arnold Networks}

\textbf{Liu et al.\ (2024)} \cite{liu2024} introduced KANs as a novel neural
network architecture inspired by the Kolmogorov-Arnold representation theorem.
Their seminal work demonstrated that KANs achieve superior accuracy and data
efficiency compared to traditional MLPs on multiple benchmark datasets. Key
identified challenges include slower training speed due to B-spline computation
complexity.

\textbf{Hou et al.\ (2024)} \cite{hou2024} conducted a comprehensive survey
reviewing KAN variants and applications across multiple domains, highlighting the
diversity of KAN architectures, theoretical foundations, and directions for
efficiency improvements.

\section{KANs for Sequential and Temporal Tasks}

\textbf{Xu, Chen, and Wang (2024)} \cite{xu2024kan} explored KANs for time
series analysis, demonstrating how the architecture bridges predictive power and
interpretability in temporal data. Their work established that KAN's learnable
univariate functions effectively capture temporal patterns --- directly relevant
to session-based recommendations.

\textbf{Genet and Inzirillo (2024)} \cite{genet2024} proposed TKAN (Temporal
Kolmogorov-Arnold Networks), specifically designed for temporal sequence modeling.
This work validated the applicability of KAN architectures to sequential
prediction tasks, providing architectural patterns that inform the design of
KAN-RNN hybrids.

\textbf{Xu, Fang, Gong et al.\ (2026)} \cite{xu2026} recently published work on
enhancing sequential recommendations using KANs, focusing on fusing global and
local interest patterns. Their research demonstrated that KAN-based architectures
can effectively model both short-term session dynamics and broader user
preferences.

\section{Research Gap}

The literature review reveals a critical gap: while KANs demonstrate superior
interpretability and parameter efficiency in various domains, and session-based
recommendation systems have evolved through RNN and Transformer architectures,
\textbf{no prior work has systematically explored KAN-based recurrent cells for
session-based recommendations} --- specifically building fully recurrent KAN cells
or stacking KAN refinement on top of GRU encoders. This project directly addresses
this gap by benchmarking three architectures under identical training conditions
on the YooChoose dataset.

% ═══════════════════════════════════════════════════════════════
% CHAPTER 3: DATASET AND PREPROCESSING
% ═══════════════════════════════════════════════════════════════
\chapter{Dataset and Preprocessing}

\section{YooChoose Dataset}

The \textbf{YooChoose} dataset originates from the ACM RecSys Challenge 2015 and
contains anonymized click-stream data from a European e-commerce platform. It is
one of the most widely used benchmarks for session-based recommendation research.
The raw dataset contains 33,003,944 interaction rows spanning 9,249,729 unique
sessions.

\subsection{Dataset Statistics}

\begin{table}[H]
\centering
\caption{YooChoose Dataset Statistics (After Preprocessing)}
\begin{tabular}{lr}
\toprule
\textbf{Property} & \textbf{Value} \\
\midrule
Raw rows (full dataset)         & 33,003,944 \\
Raw unique sessions (full)      & 9,249,729  \\
Sessions sampled                & 50,000     \\
Unique items after encoding     & 11,910     \\
Minimum session length          & 3          \\
Maximum sequence length         & 20         \\
Training samples                & 102,423    \\
Testing samples                 & 25,606     \\
Total samples                   & 128,029    \\
Train/Test split                & 80\% / 20\% \\
\bottomrule
\end{tabular}
\label{tab:dataset}
\end{table}

Each row in the raw dataset has the format:
\begin{center}
\texttt{SessionId, Timestamp, ItemId, Category}
\end{center}

\section{Preprocessing Pipeline}

The full preprocessing pipeline consists of the following steps:

\subsection{Session Filtering}

Sessions with fewer than 3 items are removed to ensure meaningful sequential
context. Sessions are trimmed to a maximum length of 20 to prevent excessive
padding. From the full 9.2M sessions, 50,000 sessions are sampled for
computational tractability on the HPC cluster.

\subsection{Item Encoding}

All raw item IDs are label-encoded to a contiguous integer range starting from 1.
Index 0 is reserved as the padding token. The vocabulary size
$|\mathcal{V}| = 11{,}910 + 1 = 11{,}911$.

\subsection{Training Sample Generation}

Each session of length $n$ is decomposed into $n-1$ training samples using a
sliding window:

\begin{equation}
\{(i_1), i_2\}, \quad \{(i_1, i_2), i_3\}, \quad \ldots, \quad \{(i_1, \ldots, i_{n-1}), i_n\}
\end{equation}

This means a session of 8 items produces 7 training samples, maximizing data
utilization. In total, 50,000 sessions produce 128,029 training/test samples.

\subsection{Padding and Batching}

Variable-length sequences within each batch are padded to the length of the
longest sequence in that batch using \texttt{pad\_sequence} with
\texttt{padding\_value=0}. Sequences longer than \texttt{MAX\_SEQ\_LEN=20} are
truncated from the left, keeping the most recent items.

\subsection{Train/Test Split}

Samples are randomly shuffled and split 80/20 into training (102,423 samples)
and test sets (25,606 samples).

\section{Dataset Visualizations}

The following figures provide an exploratory overview of the YooChoose dataset
at both the raw and preprocessed stages, helping characterise the data
distributions that drive all modelling decisions.

% ----- Row 1: Session length + Sessions per item ----- %
\begin{figure}[H]
  \centering
  \begin{subfigure}[t]{0.48\textwidth}
    \centering
    \includegraphics[width=\textwidth, trim={0 740px 1090px 0}, clip]{eda_figure.png}
    \caption{Distribution of items per session (raw). The histogram reveals a
             heavily right-skewed distribution with mean 3.57 and median 2,
             confirming the prevalence of short sessions.}
    \label{fig:dist_items_session_raw}
  \end{subfigure}
  \hfill
  \begin{subfigure}[t]{0.48\textwidth}
    \centering
    \includegraphics[width=\textwidth, trim={530px 740px 520px 0}, clip]{eda_figure.png}
    \caption{Distribution of sessions per item (raw). Most items appear in
             very few sessions (mean 625.80, median 22), reflecting the
             long-tail nature of item popularity.}
    \label{fig:dist_sessions_item_raw}
  \end{subfigure}
  \caption{Raw YooChoose interaction distributions before any preprocessing.}
  \label{fig:raw_distributions}
\end{figure}

% ----- Top-20 items (full width) ----- %
\begin{figure}[H]
  \centering
  \includegraphics[width=0.75\textwidth, trim={1090px 740px 0 0}, clip]{eda_figure.png}
  \caption{Top 20 most popular items. The strongly skewed popularity distribution,
           where the single top item (\texttt{643078800}) accumulates over
           150{,}000 sessions, motivates the item-frequency filtering step.}
  \label{fig:top20_items}
\end{figure}

% ----- Train/Test pie + sequence length distributions ----- %
\begin{figure}[H]
  \centering
  \begin{subfigure}[t]{0.28\textwidth}
    \centering
    \includegraphics[width=\textwidth, trim={100px 400px 1190px 340px}, clip]{eda_figure.png}
    \caption{Train\,/\,Test split (80\,\%\,/\,20\,\%).}
    \label{fig:train_test_split}
  \end{subfigure}
  \hfill
  \begin{subfigure}[t]{0.34\textwidth}
    \centering
    \includegraphics[width=\textwidth, trim={500px 400px 550px 360px}, clip]{eda_figure.png}
    \caption{Training sequence length distribution (mean~4.85).}
    \label{fig:train_seq_dist}
  \end{subfigure}
  \hfill
  \begin{subfigure}[t]{0.34\textwidth}
    \centering
    \includegraphics[width=\textwidth, trim={1120px 400px 0 340px}, clip]{eda_figure.png}
    \caption{Test sequence length distribution (mean~4.80).}
    \label{fig:test_seq_dist}
  \end{subfigure}
  \caption{Train/Test split and sequence-length distributions after the
           sliding-window preprocessing step. Both splits share nearly
           identical length profiles, confirming a representative random split.}
  \label{fig:split_seqlength}
\end{figure}

% ----- Cumulative distribution + dataset size comparison ----- %
\begin{figure}[H]
  \centering
  \begin{subfigure}[t]{0.48\textwidth}
    \centering
    \includegraphics[width=\textwidth, trim={0 0 1050px 680px}, clip]{eda_figure.png}
    \caption{Cumulative distribution of sessions vs.\ fraction of items.
             80\,\% of all sessions are covered by just 42{,}191 items,
             quantifying the heavy concentration of user attention.}
    \label{fig:cumulative_sessions}
  \end{subfigure}
  \hfill
  \begin{subfigure}[t]{0.48\textwidth}
    \centering
    \includegraphics[width=\textwidth, trim={530px 0 520px 680px}, clip]{eda_figure.png}
    \caption{Dataset size comparison (unique items, sessions, raw interactions,
             and preprocessed train/test sequences) highlighting the scale
             reduction achieved by the 50{,}000-session sampling strategy.}
    \label{fig:dataset_size}
  \end{subfigure}
  \caption{Item popularity concentration and final dataset size overview for
           the YooChoose 1/64 split used in this study.}
  \label{fig:cumulative_and_size}
\end{figure}

% ═══════════════════════════════════════════════════════════════
% CHAPTER 4: MODEL ARCHITECTURES
% ═══════════════════════════════════════════════════════════════
\chapter{Model Architectures}

All three models share a common skeleton:

\begin{equation}
\text{Input} \xrightarrow{\text{Embedding}(N,\,128)} \xrightarrow{\text{Sequence Encoder}} \xrightarrow{\text{Last hidden state}} \xrightarrow{\text{Dropout}(0.3)} \xrightarrow{\text{FC}} \text{Item scores}
\end{equation}

The models differ only in the \textbf{Sequence Encoder} component.

\section{Architecture Overview}

\begin{table}[H]
\centering
\caption{Model Architecture Summary}
\begin{tabular}{p{2.2cm} p{7.5cm} r r}
\toprule
\textbf{Model} & \textbf{Encoder Stack} & \textbf{Params} & \textbf{Time (s)} \\
\midrule
RNN Baseline
  & Embedding(11910, 128) $\to$ GRU(128, 256, layers=2)
    $\to$ LayerNorm $\to$ Dropout(0.3) $\to$ FC(256, 11910)
  & 5,276,805 & 145.8 \\[4pt]
Pure KAN
  & Embedding(11910, 128) $\to$ KANRecurrent$\times$2(128/256, 256)
    $\to$ LayerNorm $\to$ Dropout(0.3) $\to$ FC(256, 11910)
  & 7,569,029 & 1569.8 \\[4pt]
Hybrid
  & Embedding(11910, 128) $\to$ GRU(128, 256)
    $\to$ KANRecurrent(256, 256)
    $\to$ LayerNorm $\to$ Dropout(0.3) $\to$ FC(256, 11910)
  & 6,586,757 & 981.6 \\
\bottomrule
\end{tabular}
\label{tab:arch_overview}
\end{table}

\section{KAN Layer --- Core Building Block}

All KAN-based architectures use a shared \textbf{KANLinear} layer implemented as
B-spline basis functions. For an input $\mathbf{x} \in \mathbb{R}^{d_{in}}$, the
output is:

\begin{equation}
\text{KANLinear}(\mathbf{x}) =
  \underbrace{\sum_{d=1}^{d_{in}} \sum_{k=1}^{K} c_{o,d,k} \cdot B_k(x_d)}_{\text{spline path}}
  + \underbrace{\mathbf{W}\mathbf{x} + \mathbf{b}}_{\text{residual linear path}}
\end{equation}

where $B_k$ are B-spline basis functions of order 3 with grid size 8,
$c_{o,d,k}$ are learnable spline coefficients, and the residual linear path
ensures stable gradient flow.

The B-spline basis is computed as:
\begin{equation}
B_k^{(0)}(x) = \mathbb{1}[t_k \leq x < t_{k+1}]
\end{equation}
\begin{equation}
B_k^{(p)}(x) = \frac{x - t_k}{t_{k+p} - t_k} B_k^{(p-1)}(x)
             + \frac{t_{k+p+1} - x}{t_{k+p+1} - t_{k+1}} B_{k+1}^{(p-1)}(x)
\end{equation}

where $t_k$ are knot positions uniformly distributed on $[-1, 1]$.

\section{Model 1: RNN Baseline}

The baseline model uses a standard 2-layer GRU:

\begin{equation}
z_t = \sigma(W_z [x_t, h_{t-1}] + b_z)
\end{equation}
\begin{equation}
r_t = \sigma(W_r [x_t, h_{t-1}] + b_r)
\end{equation}
\begin{equation}
\tilde{h}_t = \tanh(W_n [x_t + r_t \odot h_{t-1}] + b_n)
\end{equation}
\begin{equation}
h_t = (1 - z_t) \odot h_{t-1} + z_t \odot \tilde{h}_t
\end{equation}

\noindent\textbf{Full stack:}
$\text{Embedding}(11910, 128) \to \text{GRU}(128, 256, \text{layers}=2)
 \to \text{LayerNorm} \to \text{Dropout}(0.3) \to \text{FC}(256, 11910)$

\noindent\textbf{Trainable parameters:} 5,276,805

\section{Model 2: Pure Recurrent KAN}

The GRU is \textbf{entirely replaced} by a custom recurrent KAN layer. At each
time step $t$:

\begin{equation}
\text{combined}_t = [\mathbf{x}_t \; \| \; \mathbf{h}_{t-1}] \in \mathbb{R}^{d_{emb} + d_{hid}}
\end{equation}
\begin{equation}
\mathbf{h}_t = \tanh\left(\text{LayerNorm}\left(\text{KAN}(\text{combined}_t)\right)\right)
\end{equation}

The hidden state $\mathbf{h}_t$ is fed back as input at the next time step,
creating true recurrence through the KAN layer. Two such layers are stacked,
where the output of layer $\ell$ becomes the input of layer $\ell+1$.
The KAN layers use grid size 8 and spline order 3.

\noindent\textbf{Full stack:}
$\text{Embedding}(11910, 128) \to \text{KANRecurrent}^{(1)}(128\to256) \to
 \text{KANRecurrent}^{(2)}(256\to256) \to \text{LayerNorm} \to \text{Dropout}(0.3) \to \text{FC}(256, 11910)$

\noindent\textbf{Trainable parameters:} 7,569,029

\section{Model 3: Hybrid (GRU + Recurrent KAN)}

Two complete recurrent systems are stacked in sequence:

\begin{enumerate}
    \item \textbf{GRU Layer}: Processes raw embeddings, captures basic sequential
          order and item co-occurrence patterns.
    \item \textbf{KAN Recurrent Layer}: Processes GRU outputs, applies non-linear
          learnable B-spline transformations over the encoded sequence.
\end{enumerate}

At each step in the KAN layer,
$\mathbf{h}^{KAN}_t = \text{KAN}([\mathbf{h}^{GRU}_t \; \| \; \mathbf{h}^{KAN}_{t-1}])$,
giving the KAN layer its own separate memory that evolves over the GRU's outputs.

\noindent\textbf{Full stack:}
$\text{Embedding}(11910, 128) \to \text{GRU}(128, 256) \to
 \text{KANRecurrent}(256, 256) \to \text{LayerNorm} \to \text{Dropout}(0.3) \to \text{FC}(256, 11910)$

\noindent\textbf{Trainable parameters:} 6,586,757

\section{Architecture Comparison Summary}

\begin{table}[H]
\centering
\caption{Detailed Architecture Comparison}
\begin{tabular}{p{2.2cm} p{2.8cm} p{2.8cm} p{2.5cm} r}
\toprule
\textbf{Model}   & \textbf{Recurrence}         & \textbf{KAN Location} & \textbf{Activation}    & \textbf{Params}  \\
\midrule
RNN Baseline
  & 2-layer GRU
  & None
  & Fixed sigmoid / tanh
  & 5,276,805 \\[4pt]
Pure KAN
  & 2$\times$ KAN cell \newline $h_t{=}\text{KAN}([x_t,h_{t-1}])$
  & Full recurrent cell
  & Learnable B-spline
  & 7,569,029 \\[4pt]
Hybrid
  & GRU then KAN stacked
  & After GRU encoder
  & GRU gates + B-spline
  & 6,586,757 \\
\bottomrule
\end{tabular}
\label{tab:arch_summary}
\end{table}

% ═══════════════════════════════════════════════════════════════
% CHAPTER 5: EXPERIMENTAL SETUP
% ═══════════════════════════════════════════════════════════════
\chapter{Experimental Setup}

\section{Hyperparameters}

All three models are trained with identical hyperparameters to ensure a fair
comparison:

\begin{table}[H]
\centering
\caption{Hyperparameter Configuration (Actual Training Run)}
\begin{tabular}{ll}
\toprule
\textbf{Hyperparameter}          & \textbf{Value}        \\
\midrule
Embedding Dimension              & 128                   \\
Hidden Dimension                 & 256                   \\
KAN Grid Size                    & 8                     \\
KAN Spline Order                 & 3                     \\
Dropout Rate                     & 0.3                   \\
Batch Size                       & 512                   \\
Epochs                           & 20                    \\
Learning Rate                    & $5 \times 10^{-4}$    \\
Weight Decay                     & $1 \times 10^{-5}$    \\
Warmup Epochs                    & 3                     \\
Optimizer                        & AdamW                 \\
LR Schedule                      & Warmup + Cosine Decay \\
Random Seed                      & 42                    \\
Device                           & CUDA (GPU)            \\
\bottomrule
\end{tabular}
\label{tab:hyperparams}
\end{table}

\section{Training Procedure}

\subsection{Loss Function}

Cross-entropy loss is used for next-item classification. For a session $S$ with
target item $i_n$ and item vocabulary $\mathcal{V}$:

\begin{equation}
\mathcal{L} = -\log p(i_n \mid S)
\end{equation}

\subsection{Learning Rate Schedule}

A linear warmup followed by cosine decay schedule is used:

\begin{equation}
\text{lr}(t) =
\begin{cases}
\text{lr}_{max} \cdot \dfrac{t}{T_{warmup}} & t < T_{warmup} \\[8pt]
\text{lr}_{max} \cdot \dfrac{1}{2}\!\left(1 + \cos\!\left(\pi \cdot
  \dfrac{t - T_{warmup}}{T_{total} - T_{warmup}}\right)\right) & t \geq T_{warmup}
\end{cases}
\end{equation}

The RNN model converges quickly and benefits from warmup during epochs 1--3.
KAN models converge more slowly due to the added spline complexity, reaching
near-plateau performance by epoch 16.

\subsection{Weight Initialization}

\begin{itemize}
    \item Embedding layer: Normal distribution ($\mu=0$, $\sigma=0.01$)
    \item FC layers: Xavier uniform initialization
    \item Bias terms: Zero initialization
\end{itemize}

\subsection{Gradient Clipping}

Gradient norm is clipped to 1.0 to prevent exploding gradients.

\section{Evaluation Metrics}

The following metrics are computed on the held-out test set after each epoch.
All ranking metrics treat next-item prediction as a ranking task over all
11,910 items.

\begin{table}[H]
\centering
\caption{Evaluation Metrics Definitions}
\begin{tabular}{p{2.5cm} p{9cm}}
\toprule
\textbf{Metric}   & \textbf{Definition} \\
\midrule
Loss (CrossEntropy)
  & Categorical cross-entropy on validation set. Lower $(\downarrow)$ is better. \\[4pt]
Accuracy (Top-1)
  & Fraction of test samples where the target item is the top-1 predicted item. Higher $(\uparrow)$ is better. \\[4pt]
HR@10
  & Hit Rate at 10. Fraction of sessions where the target item appears in the top-10 ranked predictions. Higher $(\uparrow)$ is better. \\[4pt]
HR@20
  & Hit Rate at 20. As above but within the top-20 predictions. Higher $(\uparrow)$ is better. \\[4pt]
NDCG@10
  & Normalized Discounted Cumulative Gain at 10. Rewards correct predictions at higher rank positions. Higher $(\uparrow)$ is better. \\[4pt]
NDCG@20
  & As above within the top-20 predictions. Higher $(\uparrow)$ is better. \\[4pt]
Recall@20
  & For single-target prediction, equivalent to HR@20. Higher $(\uparrow)$ is better. \\[4pt]
MRR@20
  & Mean Reciprocal Rank at 20. Average of $1/\text{rank}(i_n)$ for sessions where the target appears in top-20. Higher $(\uparrow)$ is better. \\
\bottomrule
\end{tabular}
\label{tab:metrics_def}
\end{table}

\subsection{Formal Definitions}

\begin{equation}
\text{Acc@1} = \frac{1}{|T|} \sum_{s \in T} \mathbb{1}[\hat{i}_1(s) = i_n(s)]
\end{equation}

\begin{equation}
\text{HR@K} = \frac{1}{|T|} \sum_{s \in T} \mathbb{1}[i_n(s) \in \text{Top-K}(s)]
\end{equation}

\begin{equation}
\text{NDCG@K} = \frac{1}{|T|} \sum_{s \in T}
  \frac{\mathbb{1}[i_n(s) \in \text{Top-K}(s)]}{\log_2(\text{rank}(i_n(s)) + 1)}
\end{equation}

\begin{equation}
\text{MRR@K} = \frac{1}{|T|} \sum_{s \in T}
  \frac{\mathbb{1}[i_n(s) \in \text{Top-K}(s)]}{\text{rank}(i_n(s))}
\end{equation}

\section{Implementation Details}

\begin{itemize}
    \item \textbf{Framework}: PyTorch 2.0+
    \item \textbf{Hardware}: NVIDIA GPU (HPC cluster, CUDA)
    \item \textbf{Reproducibility}: Fixed random seed (42) for all experiments
\end{itemize}

The codebase is modular with separate files for each model:

\begin{table}[H]
\centering
\caption{Project File Structure}
\begin{tabular}{ll}
\toprule
\textbf{File}              & \textbf{Role}                                  \\
\midrule
\texttt{config.py}         & All hyperparameters                            \\
\texttt{data\_loader.py}   & YooChoose loading and preprocessing            \\
\texttt{kan\_layer.py}     & B-spline KAN linear layer                      \\
\texttt{kan\_recurrent.py} & Recurrent KAN cell and layer                   \\
\texttt{rnn\_model.py}     & Model 1: GRU baseline                          \\
\texttt{kan\_model.py}     & Model 2: Pure recurrent KAN                    \\
\texttt{hybrid\_model.py}  & Model 3: GRU + KAN stacked                     \\
\texttt{trainer.py}        & Shared training and evaluation loop            \\
\texttt{metrics.py}        & HR@10, HR@20, NDCG@10, NDCG@20, MRR           \\
\texttt{compare.py}        & Main script: trains all models, plots, table   \\
\bottomrule
\end{tabular}
\label{tab:files}
\end{table}

% ═══════════════════════════════════════════════════════════════
% CHAPTER 6: RESULTS AND DISCUSSION
% ═══════════════════════════════════════════════════════════════
\chapter{Results and Discussion}

% ---------------------------------------------------------------
\section{Epoch-by-Epoch Training Curves}
% ---------------------------------------------------------------

This chapter presents the quantitative evaluation of all three architectures on
the YooChoose test set. Tables~\ref{tab:rnn_log}--\ref{tab:hybrid_log} summarize
the per-epoch training and validation metrics over 20 epochs. The corresponding
learning curves are presented in Figures~\ref{fig:loss_curves}--\ref{fig:rank_curves}.

\subsection{RNN Baseline Training Log}

\begin{table}[H]
\centering
\caption{RNN Baseline --- Per-Epoch Metrics (Total time: 145.8 s)}
\small
\begin{tabular}{r r r r r r r r}
\toprule
\textbf{Ep} & \textbf{TrLoss} & \textbf{VlLoss} & \textbf{Acc} & \textbf{HR@10} & \textbf{HR@20} & \textbf{NDCG@10} & \textbf{MRR} \\
\midrule
 1  & 8.7991 & 8.1339 & 0.0075 & 0.0614 & 0.1009 & 0.0305 & 0.0238 \\
 2  & 7.9848 & 7.5568 & 0.0518 & 0.1802 & 0.2343 & 0.1080 & 0.0896 \\
 3  & 7.0007 & 6.4887 & 0.1271 & 0.3801 & 0.4571 & 0.2444 & 0.2073 \\
 4  & 6.0661 & 5.9280 & 0.1723 & 0.4802 & 0.5567 & 0.3183 & 0.2728 \\
 5  & 5.4881 & 5.6646 & 0.1984 & 0.5273 & 0.6002 & 0.3543 & 0.3049 \\
 6  & 5.1113 & 5.5369 & 0.2063 & 0.5502 & 0.6256 & 0.3707 & 0.3193 \\
 7  & 4.8530 & 5.4860 & 0.2169 & 0.5663 & 0.6363 & 0.3828 & 0.3299 \\
 8  & 4.6511 & 5.4609 & 0.2166 & 0.5733 & 0.6482 & 0.3866 & 0.3330 \\
 9  & 4.4963 & 5.4649 & 0.2204 & 0.5768 & 0.6508 & 0.3899 & 0.3362 \\
10  & 4.3730 & 5.4755 & 0.2202 & 0.5787 & 0.6540 & 0.3906 & 0.3366 \\
11  & 4.2724 & 5.4963 & 0.2212 & 0.5813 & 0.6551 & 0.3925 & 0.3381 \\
12  & 4.1906 & 5.5198 & 0.2219 & 0.5803 & 0.6562 & 0.3925 & 0.3386 \\
13  & 4.1206 & 5.5329 & 0.2250 & 0.5820 & 0.6543 & 0.3942 & 0.3401 \\
14  & 4.0634 & 5.5513 & 0.2230 & 0.5810 & 0.6540 & 0.3932 & 0.3391 \\
15  & 4.0187 & 5.5611 & 0.2214 & 0.5812 & 0.6550 & 0.3926 & 0.3383 \\
16  & 3.9825 & 5.5708 & 0.2221 & 0.5814 & 0.6546 & 0.3930 & 0.3387 \\
17  & 3.9530 & 5.5758 & 0.2251 & 0.5811 & 0.6545 & 0.3941 & 0.3403 \\
18  & 3.9322 & 5.5788 & 0.2253 & 0.5809 & 0.6545 & 0.3942 & 0.3405 \\
19  & 3.9244 & 5.5823 & 0.2255 & 0.5811 & 0.6543 & 0.3944 & 0.3406 \\
\textbf{20} & \textbf{3.9221} & \textbf{5.5849} & \textbf{0.2253} & \textbf{0.5808} & \textbf{0.6542} & \textbf{0.3942} & \textbf{0.3405} \\
\bottomrule
\end{tabular}
\label{tab:rnn_log}
\end{table}

\subsection{Pure KAN Training Log}

\begin{table}[H]
\centering
\caption{Pure KAN --- Per-Epoch Metrics (Total time: 1569.8 s)}
\small
\begin{tabular}{r r r r r r r r}
\toprule
\textbf{Ep} & \textbf{TrLoss} & \textbf{VlLoss} & \textbf{Acc} & \textbf{HR@10} & \textbf{HR@20} & \textbf{NDCG@10} & \textbf{MRR} \\
\midrule
 1  & 8.8652 & 8.1575 & 0.0088 & 0.0616 & 0.0981 & 0.0295 & 0.0225 \\
 2  & 8.1204 & 8.1154 & 0.0110 & 0.0634 & 0.0997 & 0.0322 & 0.0254 \\
 3  & 7.8649 & 7.5769 & 0.0249 & 0.1312 & 0.1869 & 0.0707 & 0.0562 \\
 4  & 7.2501 & 7.0365 & 0.0523 & 0.2401 & 0.3225 & 0.1352 & 0.1086 \\
 5  & 6.7105 & 6.6397 & 0.0918 & 0.3290 & 0.4086 & 0.1990 & 0.1644 \\
 6  & 6.2730 & 6.3846 & 0.1160 & 0.3867 & 0.4650 & 0.2402 & 0.2001 \\
 7  & 5.9322 & 6.2095 & 0.1391 & 0.4311 & 0.5027 & 0.2745 & 0.2306 \\
 8  & 5.6675 & 6.0776 & 0.1575 & 0.4555 & 0.5295 & 0.2967 & 0.2522 \\
 9  & 5.4419 & 6.0083 & 0.1651 & 0.4717 & 0.5458 & 0.3091 & 0.2632 \\
10  & 5.2529 & 5.9607 & 0.1751 & 0.4878 & 0.5585 & 0.3227 & 0.2758 \\
11  & 5.0952 & 5.9126 & 0.1808 & 0.4984 & 0.5692 & 0.3317 & 0.2842 \\
12  & 4.9622 & 5.8991 & 0.1889 & 0.5037 & 0.5746 & 0.3377 & 0.2905 \\
13  & 4.8561 & 5.8900 & 0.1930 & 0.5088 & 0.5777 & 0.3416 & 0.2939 \\
14  & 4.7610 & 5.8849 & 0.1955 & 0.5103 & 0.5797 & 0.3446 & 0.2973 \\
15  & 4.6886 & 5.8773 & 0.1999 & 0.5140 & 0.5821 & 0.3486 & 0.3013 \\
16  & 4.6326 & 5.8670 & 0.2006 & 0.5161 & 0.5843 & 0.3495 & 0.3019 \\
17  & 4.5915 & 5.8671 & 0.1981 & 0.5158 & 0.5854 & 0.3487 & 0.3010 \\
18  & 4.5660 & 5.8654 & 0.2006 & 0.5166 & 0.5850 & 0.3503 & 0.3027 \\
19  & 4.5501 & 5.8653 & 0.2012 & 0.5170 & 0.5856 & 0.3511 & 0.3037 \\
\textbf{20} & \textbf{4.5414} & \textbf{5.8678} & \textbf{0.2020} & \textbf{0.5174} & \textbf{0.5850} & \textbf{0.3511} & \textbf{0.3036} \\
\bottomrule
\end{tabular}
\label{tab:kan_log}
\end{table}

\subsection{Hybrid Model Training Log}

\begin{table}[H]
\centering
\caption{Hybrid (GRU + KAN) --- Per-Epoch Metrics (Total time: 981.6 s)}
\small
\begin{tabular}{r r r r r r r r}
\toprule
\textbf{Ep} & \textbf{TrLoss} & \textbf{VlLoss} & \textbf{Acc} & \textbf{HR@10} & \textbf{HR@20} & \textbf{NDCG@10} & \textbf{MRR} \\
\midrule
 1  & 8.8390 & 8.1555 & 0.0081 & 0.0616 & 0.0993 & 0.0293 & 0.0222 \\
 2  & 8.0955 & 7.9467 & 0.0092 & 0.0635 & 0.1020 & 0.0311 & 0.0241 \\
 3  & 7.5964 & 7.1895 & 0.0529 & 0.2246 & 0.2963 & 0.1281 & 0.1034 \\
 4  & 6.7455 & 6.5120 & 0.1106 & 0.3573 & 0.4309 & 0.2251 & 0.1889 \\
 5  & 6.1028 & 6.1547 & 0.1501 & 0.4286 & 0.4991 & 0.2804 & 0.2389 \\
 6  & 5.6772 & 5.9623 & 0.1732 & 0.4686 & 0.5437 & 0.3119 & 0.2681 \\
 7  & 5.3547 & 5.8428 & 0.1860 & 0.4944 & 0.5691 & 0.3316 & 0.2856 \\
 8  & 5.1008 & 5.7726 & 0.1962 & 0.5115 & 0.5814 & 0.3456 & 0.2983 \\
 9  & 4.9018 & 5.7400 & 0.2045 & 0.5210 & 0.5903 & 0.3547 & 0.3073 \\
10  & 4.7219 & 5.7234 & 0.2059 & 0.5284 & 0.5970 & 0.3592 & 0.3108 \\
11  & 4.5789 & 5.7211 & 0.2098 & 0.5318 & 0.6048 & 0.3626 & 0.3144 \\
12  & 4.4557 & 5.7161 & 0.2139 & 0.5358 & 0.6024 & 0.3669 & 0.3184 \\
13  & 4.3579 & 5.7238 & 0.2137 & 0.5359 & 0.6040 & 0.3667 & 0.3182 \\
14  & 4.2703 & 5.7247 & 0.2132 & 0.5368 & 0.6055 & 0.3667 & 0.3180 \\
15  & 4.2056 & 5.7311 & 0.2155 & 0.5372 & 0.6063 & 0.3679 & 0.3194 \\
16  & 4.1578 & 5.7354 & 0.2170 & 0.5381 & 0.6061 & 0.3687 & 0.3202 \\
17  & 4.1179 & 5.7358 & 0.2156 & 0.5389 & 0.6069 & 0.3684 & 0.3195 \\
18  & 4.0946 & 5.7361 & 0.2157 & 0.5387 & 0.6068 & 0.3685 & 0.3197 \\
19  & 4.0856 & 5.7385 & 0.2165 & 0.5386 & 0.6063 & 0.3688 & 0.3201 \\
\textbf{20} & \textbf{4.0736} & \textbf{5.7374} & \textbf{0.2169} & \textbf{0.5386} & \textbf{0.6071} & \textbf{0.3690} & \textbf{0.3205} \\
\bottomrule
\end{tabular}
\label{tab:hybrid_log}
\end{table}

% ---------------------------------------------------------------
\section{Training Curve Comparisons}
% ---------------------------------------------------------------

Figure~\ref{fig:loss_curves} presents the training and validation loss alongside
top-1 accuracy. Figures~\ref{fig:hit_curves} and~\ref{fig:rank_curves} show the
retrieval and ranking metrics respectively.

% ----- Figure: Loss + Accuracy (top row of training_curves.png) ----- %
\begin{figure}[H]
  \centering
  \begin{subfigure}[t]{0.32\textwidth}
    \centering
    \includegraphics[width=\textwidth, trim={0 660px 1080px 0}, clip]{training_curves.png}
    \caption{Training loss over 20 epochs. The RNN converges fastest,
             reaching near-plateau by epoch~10. KAN and Hybrid decrease
             monotonically, suggesting further gains beyond 20 epochs.}
    \label{fig:train_loss}
  \end{subfigure}
  \hfill
  \begin{subfigure}[t]{0.32\textwidth}
    \centering
    \includegraphics[width=\textwidth, trim={540px 660px 540px 0}, clip]{training_curves.png}
    \caption{Validation loss over 20 epochs. The RNN achieves the lowest
             minimum ($\approx$5.46 at epoch~8) before mild overfitting;
             KAN and Hybrid continue improving.}
    \label{fig:val_loss}
  \end{subfigure}
  \hfill
  \begin{subfigure}[t]{0.32\textwidth}
    \centering
    \includegraphics[width=\textwidth, trim={1080px 660px 0 0}, clip]{training_curves.png}
    \caption{Top-1 validation accuracy over 20 epochs. The RNN peaks at
             0.2255 (epoch~19); KAN and Hybrid converge more gradually.}
    \label{fig:val_acc}
  \end{subfigure}
  \caption{Loss and accuracy learning curves for RNN (red), KAN (green),
           and Hybrid (blue) models over 20 training epochs.}
  \label{fig:loss_curves}
\end{figure}

% ----- Figure: HR@10, HR@20, NDCG@10 (middle row) ----- %
\begin{figure}[H]
  \centering
  \begin{subfigure}[t]{0.32\textwidth}
    \centering
    \includegraphics[width=\textwidth, trim={0 340px 1080px 360px}, clip]{training_curves.png}
    \caption{HR@10 over 20 epochs. The RNN achieves the highest HR@10
             (0.5808); Hybrid reaches 0.5386 and Pure KAN 0.5174.}
    \label{fig:hr10}
  \end{subfigure}
  \hfill
  \begin{subfigure}[t]{0.32\textwidth}
    \centering
    \includegraphics[width=\textwidth, trim={540px 340px 540px 360px}, clip]{training_curves.png}
    \caption{HR@20 over 20 epochs. Consistent ordering:
             RNN (0.6542) $>$ Hybrid (0.6071) $>$ KAN (0.5850).}
    \label{fig:hr20}
  \end{subfigure}
  \hfill
  \begin{subfigure}[t]{0.32\textwidth}
    \centering
    \includegraphics[width=\textwidth, trim={1080px 340px 0 360px}, clip]{training_curves.png}
    \caption{NDCG@10 over 20 epochs. The RNN's ranking quality is superior
             (0.3942), placing relevant items higher in the top-10 list.}
    \label{fig:ndcg10}
  \end{subfigure}
  \caption{Hit-rate and NDCG@10 learning curves. The RNN establishes clear
           superiority from epoch~5 onward; the Hybrid sits between RNN and
           Pure KAN on all metrics.}
  \label{fig:hit_curves}
\end{figure}

% ----- Figure: NDCG@20, Recall@10, MRR (bottom row) ----- %
\begin{figure}[H]
  \centering
  \begin{subfigure}[t]{0.32\textwidth}
    \centering
    \includegraphics[width=\textwidth, trim={0 0 1080px 710px}, clip]{training_curves.png}
    \caption{NDCG@20 over 20 epochs. The RNN advantage is maintained at the
             wider rank cut-off (0.4128 vs.\ 0.3864 Hybrid, 0.3683 KAN).}
    \label{fig:ndcg20}
  \end{subfigure}
  \hfill
  \begin{subfigure}[t]{0.32\textwidth}
    \centering
    \includegraphics[width=\textwidth, trim={540px 0 540px 710px}, clip]{training_curves.png}
    \caption{Recall@10 over 20 epochs. Values mirror HR@10 in this
             single-ground-truth setting, confirming consistency.}
    \label{fig:recall10}
  \end{subfigure}
  \hfill
  \begin{subfigure}[t]{0.32\textwidth}
    \centering
    \includegraphics[width=\textwidth, trim={1080px 0 0 710px}, clip]{training_curves.png}
    \caption{MRR@20 over 20 epochs. The RNN (0.3405) consistently ranks the
             correct item higher than Hybrid (0.3205) or KAN (0.3036).}
    \label{fig:mrr}
  \end{subfigure}
  \caption{NDCG@20, Recall@10, and MRR@20 learning curves confirming the
           consistent performance ordering across all ranking metrics.}
  \label{fig:rank_curves}
\end{figure}

% ---------------------------------------------------------------
\section{Final Quantitative Comparison}
% ---------------------------------------------------------------

Table~\ref{tab:results} presents the definitive comparison of all three models
at epoch 20 on the YooChoose test set, along with the relative improvement or
degradation of KAN-based models versus the RNN baseline.

\begin{table}[H]
\centering
\caption{Final Model Comparison --- All Metrics on YooChoose Test Set (Epoch 20)}
\begin{tabular}{l r r r r r}
\toprule
\textbf{Metric} & \textbf{RNN} & \textbf{KAN} & \textbf{Hybrid}
                & \textbf{KAN vs RNN} & \textbf{Hybrid vs RNN} \\
\midrule
Val Loss $\downarrow$  & 5.5849 & 5.8678 & 5.7374 & $-$0.2829 & $-$0.1525 \\
Accuracy $\uparrow$    & 0.2253 & 0.2020 & 0.2169 & $-$0.0233 & $-$0.0084 \\
HR@10 $\uparrow$       & 0.5808 & 0.5174 & 0.5386 & $-$0.0634 & $-$0.0422 \\
HR@20 $\uparrow$       & 0.6542 & 0.5850 & 0.6071 & $-$0.0692 & $-$0.0471 \\
NDCG@10 $\uparrow$     & 0.3942 & 0.3511 & 0.3690 & $-$0.0431 & $-$0.0252 \\
NDCG@20 $\uparrow$     & 0.4128 & 0.3683 & 0.3864 & $-$0.0445 & $-$0.0264 \\
Recall@20 $\uparrow$   & 0.5808 & 0.5174 & 0.5386 & $-$0.0634 & $-$0.0422 \\
MRR@20 $\uparrow$      & 0.3405 & 0.3036 & 0.3205 & $-$0.0369 & $-$0.0200 \\
\midrule
Time (s)               & \textbf{145.8}   & 1569.8   & 981.6    & $+$1424.0 s & $+$835.8 s \\
Params                 & \textbf{5,276,805} & 7,569,029 & 6,586,757 & $+$2.29M & $+$1.31M \\
\midrule
\multicolumn{3}{l}{\textbf{Winner (metrics)} $\uparrow$} & \multicolumn{3}{l}{\textbf{RNN (7/7 metrics)}} \\
\bottomrule
\multicolumn{6}{l}{\small $\uparrow$ = higher is better; $\downarrow$ = lower is better. Negative deltas indicate KAN/Hybrid underperformance.} \\
\end{tabular}
\label{tab:results}
\end{table}

% ---------------------------------------------------------------
\section{Architecture Analysis}
% ---------------------------------------------------------------

\subsection{RNN Baseline}

The standard 2-layer GRU baseline achieves the best performance across all
evaluation metrics while being by far the fastest to train (145.8 s, 10.8$\times$
faster than Pure KAN). The gating mechanism provides effective memory management
over the session sequence, and the 2-layer depth gives sufficient representational
capacity for the 11,910-item vocabulary. Training loss decreases steadily, with
validation loss plateauing around epoch 8--9 (minimum at 5.4609) before slowly
increasing due to mild overfitting.

\subsection{Pure Recurrent KAN}

Replacing the GRU entirely with recurrent KAN cells introduces fully learnable
B-spline transformation curves but at substantial computational cost:
10.8$\times$ longer training time and 43\% more parameters than the RNN baseline.
The model converges significantly more slowly --- only reaching HR@10 = 0.50 by
epoch 14, compared to the RNN achieving HR@10 = 0.48 by epoch 4. The final
performance gap (HR@10: 0.5174 vs.\ 0.5808) suggests that pure spline recurrence
without the gating structure of GRU struggles to resolve the long-range dependency
and vanishing gradient problems that GRU gates directly address. However, the KAN
model does eventually converge, validating that recurrent KAN cells are viable
for session-based recommendation.

\subsection{Hybrid Model}

The Hybrid model represents an intermediate point: it outperforms Pure KAN on all
metrics while training in roughly 67\% of the KAN training time, but still falls
short of the RNN baseline. The two-stage design --- GRU for sequential order
capture, KAN for non-linear refinement --- partially addresses the slow convergence
issue by providing the KAN layer with already-meaningful GRU-encoded representations
rather than raw embeddings. The Hybrid model shows the most consistent improvement
across epochs 5--12, suggesting the KAN refinement layer adds value when built on
top of a well-initialized sequential representation.

% ---------------------------------------------------------------
\section{Training Speed and Parameter Efficiency}
% ---------------------------------------------------------------

\begin{table}[H]
\centering
\caption{Computational Cost Comparison}
\begin{tabular}{l r r r}
\toprule
\textbf{Model} & \textbf{Params} & \textbf{Training Time (s)} & \textbf{Time per Epoch (s)} \\
\midrule
RNN Baseline & 5,276,805  &  145.8  &  7.3  \\
Pure KAN     & 7,569,029  & 1569.8  & 78.5  \\
Hybrid       & 6,586,757  &  981.6  & 49.1  \\
\bottomrule
\end{tabular}
\label{tab:compute}
\end{table}

The B-spline basis computation is the primary source of overhead in KAN layers.
Each forward pass requires evaluating $d_{in} \times (G + p)$ spline coefficients
per output neuron, where $G=8$ is the grid size and $p=3$ is the spline order.
This results in approximately $11\times$ slower per-epoch training for Pure KAN
versus RNN.

% ---------------------------------------------------------------
\section{B-Spline Interpretability}
% ---------------------------------------------------------------

One key advantage of KAN-based architectures over standard RNNs is the
interpretability of learned transformations. Each KAN layer learns a set of
univariate spline curves --- one per input dimension. These curves can be
visualized to understand:

\begin{itemize}
    \item Which item embedding dimensions contribute most to the recommendation
          decision (by examining spline coefficient magnitudes).
    \item What non-linear shape is learned for each recurrent transformation
          (e.g., monotone, oscillatory, piece-wise linear).
    \item How the Hybrid model's KAN refinement layer differs from the initial
          GRU encoding in its learned spline shapes.
\end{itemize}

This interpretability advantage is entirely absent in standard GRU models, which
expose only weight matrices with no direct insight into the learned transformation
shape.

% ═══════════════════════════════════════════════════════════════
% CHAPTER 7: CONCLUSION AND FUTURE WORK
% ═══════════════════════════════════════════════════════════════
\chapter{Conclusion and Future Work}

\section{Conclusion}

This project implemented and systematically benchmarked three architectures for
session-based recommendation on the YooChoose e-commerce dataset (50,000
sessions, 11,910 items, 128,029 training/test samples):

\begin{enumerate}
    \item \textbf{RNN Baseline} (5.27M params, 145.8 s): 2-layer GRU with fixed
          sigmoid/tanh activations. Achieves the best performance on all seven
          evaluation metrics --- HR@10 = 0.5808, NDCG@10 = 0.3942, Acc = 0.2253,
          MRR@20 = 0.3405.
    \item \textbf{Pure Recurrent KAN} (7.57M params, 1569.8 s): Full replacement
          of GRU with B-spline recurrent cells. Converges more slowly and achieves
          HR@10 = 0.5174, NDCG@10 = 0.3511, at $10.8\times$ the training cost.
    \item \textbf{Hybrid (GRU + KAN)} (6.59M params, 981.6 s): GRU encoder
          followed by KAN recurrent refinement. Intermediate performance (HR@10 =
          0.5386, NDCG@10 = 0.3690) and training cost, demonstrating that KAN
          refinement on top of GRU outputs partially closes the performance gap.
\end{enumerate}

The key finding is that the standard GRU gating mechanism --- which directly
addresses vanishing gradients and long-range dependencies through learned forget
and update gates --- provides a stronger inductive bias for session modeling than
pure B-spline recurrence within the 20-epoch training budget. KAN-based models
demonstrate the interpretability advantage of learnable spline curves but require
significantly more training time and parameters to approach comparable performance.

The theoretical contribution of this work is the \textbf{KANRecurrentCell} ---
a novel recurrent unit using B-spline parameterized learnable activation functions
in place of fixed activations, validated for the first time on a real-world
session-based recommendation benchmark.

\section{Future Work}

Several directions remain for future exploration:

\begin{itemize}
    \item \textbf{Extended training}: KAN models show consistent improvement
          through epoch 20 without clear plateauing, suggesting that training
          beyond 20 epochs may narrow the performance gap with RNN.
    \item \textbf{Spline visualization}: Generate and analyze learned B-spline
          activation patterns across layers to identify which embedding dimensions
          drive recommendation decisions.
    \item \textbf{Ablation studies}: Systematic variation of KAN grid size
          (4, 6, 8, 12) and order (2, 3, 4) to identify the optimal
          complexity-accuracy-speed trade-off.
    \item \textbf{KAN-GRU (gates integration)}: Extend the KAN integration inside
          each GRU gate (update, reset, new gate) using separate KAN layers for
          $x_t$ and $h_{t-1}$, with reduced grid size for memory efficiency.
    \item \textbf{Attention + KAN}: Combine attention mechanisms (as in NARM)
          with KAN layers for capturing both global session intent and local
          sequential patterns.
    \item \textbf{Larger dataset}: Evaluate on the full YooChoose dataset
          (9.2M sessions) and the Diginetica dataset.
    \item \textbf{KAN efficiency}: Explore efficient KAN variants (rational
          functions, Chebyshev polynomials) to reduce the $10\times$ speed
          overhead of B-spline KAN versus GRU.
    \item \textbf{Recall@20 and MRR@20 optimization}: Directly optimize
          ranking-aware loss functions to improve MRR rather than cross-entropy
          which optimizes for top-1 accuracy.
\end{itemize}

% ═══════════════════════════════════════════════════════════════
% APPENDIX
% ═══════════════════════════════════════════════════════════════
\appendix

\chapter{KAN Layer Implementation}

\begin{lstlisting}[style=pythonstyle, caption={KANLinear --- B-spline based KAN layer (kan\_layer.py)}]
class KANLinear(nn.Module):
    def __init__(self, in_dim, out_dim, grid=KAN_GRID, order=KAN_ORDER):
        super().__init__()
        n_coeff = grid + order
        self.coeff  = nn.Parameter(torch.randn(out_dim, in_dim, n_coeff) * 0.05)
        self.linear = nn.Linear(in_dim, out_dim, bias=True)
        self.scale  = nn.Parameter(torch.ones(out_dim, in_dim) * 0.1)
        self.register_buffer("knots", torch.linspace(-1, 1, grid + 1))

    def _bspline(self, x):
        k0, target = self.knots, self.grid + self.order
        xv = x.unsqueeze(-1)
        b  = ((xv >= k0[:-1]) & (xv < k0[1:])).float()
        for k in range(1, self.order + 1):
            kext  = torch.cat([k0[0].expand(k), k0, k0[-1].expand(k)])
            w     = b.size(-1)
            denom = (kext[k: k+w] - kext[:w]).clamp(min=1e-6)
            alpha = ((xv - kext[:w]) / denom).clamp(0.0, 1.0)
            b_sh  = F.pad(b[..., 1:], (0, 1))
            b     = alpha * b + (1.0 - alpha) * b_sh
            if b.size(-1) < target:
                b = F.pad(b, (0, target - b.size(-1)))
        return b

    def forward(self, x):
        xn     = torch.tanh(x).clamp(-1+1e-5, 1-1e-5)
        basis  = self._bspline(xn)
        spline = torch.einsum("bdc,odc->bo", basis,
                              self.coeff * self.scale.unsqueeze(-1))
        return spline + self.linear(x)
\end{lstlisting}

\chapter{Recurrent KAN Cell Implementation}

\begin{lstlisting}[style=pythonstyle, caption={KANRecurrentCell --- Fully recurrent KAN cell (kan\_recurrent.py)}]
class KANRecurrentCell(nn.Module):
    """
    Recurrent cell: h_t = tanh(LayerNorm(KAN([x_t || h_{t-1}])))
    """
    def __init__(self, input_dim, hidden_dim):
        super().__init__()
        self.kan  = KANLinear(input_dim + hidden_dim, hidden_dim)
        self.norm = nn.LayerNorm(hidden_dim)

    def forward(self, x_t, h_prev):
        combined = torch.cat([x_t, h_prev], dim=-1)
        return torch.tanh(self.norm(self.kan(combined)))


class KANRecurrentLayer(nn.Module):
    """
    Unrolls a KANRecurrentCell over a sequence (batch, seq_len, input_dim)
    and returns (outputs, final_hidden).
    """
    def __init__(self, input_dim, hidden_dim):
        super().__init__()
        self.cell       = KANRecurrentCell(input_dim, hidden_dim)
        self.hidden_dim = hidden_dim

    def forward(self, x, h0=None):
        B, T, _ = x.shape
        h = torch.zeros(B, self.hidden_dim, device=x.device) if h0 is None else h0
        outputs = []
        for t in range(T):
            h = self.cell(x[:, t, :], h)
            outputs.append(h.unsqueeze(1))
        return torch.cat(outputs, dim=1), h
\end{lstlisting}

% ═══════════════════════════════════════════════════════════════
% REFERENCES
% ═══════════════════════════════════════════════════════════════
\begin{thebibliography}{99}

\bibitem{hidasi2015}
Hidasi, B., Karatzoglou, A., Baltrunas, L., \& Tikk, D. (2015).
\textit{Session-based Recommendations with Recurrent Neural Networks}.
ArXiv. \url{https://arxiv.org/abs/1511.06939}

\bibitem{li2017}
Li, J., Ren, P., Chen, Z., Ren, Z., Lian, T., \& Ma, J. (2017).
\textit{Neural Attentive Session-based Recommendation}.
Proceedings of CIKM 2017, ACM, 1419--1428.

\bibitem{sun2019}
Sun, F., Liu, J., Wu, J., Pei, C., Lin, X., Ou, W., \& Jiang, P. (2019).
\textit{BERT4Rec: Sequential Recommendation with Bidirectional Encoder
Representations from Transformer}.
ArXiv. \url{https://arxiv.org/abs/1904.06690}

\bibitem{kang2018}
Kang, W.-C., \& McAuley, J. (2018).
\textit{Self-Attentive Sequential Recommendation}.
IEEE International Conference on Data Mining (ICDM), 197--206.

\bibitem{liu2024}
Liu, Z., Wang, Y., Vaidya, S., Ruehle, F., Halverson, J., Soljačić, M.,
Hou, T.~Y., \& Tegmark, M. (2024).
\textit{KAN: Kolmogorov-Arnold Networks}.
ArXiv. \url{https://arxiv.org/abs/2404.19756}

\bibitem{hou2024}
Hou, Y., Ji, T., Zhang, D., \& Stefanidis, A. (2024).
\textit{Kolmogorov-Arnold Networks: A Critical Assessment of Claims,
Performance, and Practical Viability}.
ArXiv. \url{https://arxiv.org/abs/2407.11075}

\bibitem{xu2024kan}
Xu, K., Chen, L., \& Wang, S. (2024).
\textit{Kolmogorov-Arnold Networks for Time Series: Bridging Predictive Power
and Interpretability}.
ArXiv. \url{https://arxiv.org/abs/2406.02496}

\bibitem{genet2024}
Genet, R., \& Inzirillo, H. (2024).
\textit{TKAN: Temporal Kolmogorov-Arnold Networks}.
ArXiv. \url{https://arxiv.org/abs/2405.07344}

\bibitem{xu2026}
Xu, Y., Fang, X., Gong, J. et al. (2026).
\textit{Enhancing global and local interests fusion based on
Kolmogorov-Arnold networks for sequential recommendation}.
Multimedia Systems, 32, 81.
\url{https://doi.org/10.1007/s00530-025-02142-4}

\bibitem{rendle2010}
Rendle, S., Freudenthaler, C., \& Schmidt-Thieme, L. (2010).
\textit{Factorizing personalized Markov chains for next-basket recommendation}.
Proceedings of WWW 2010, 811--820.

\end{thebibliography}

\end{document}
